\section{Memory Usage Analysis}
\label{sec:memory_usage_analysis}

The memory plot is drawn on a logarithmic scale, so it should not be read like a
normal line chart. A small-looking rise near \(10\) KB is not the same as a
small-looking rise near \(10{,}000\) KB. The useful question is usually not
"which line is a little higher?", but "has this solver moved into a different
order of magnitude?" Table~\ref{tab:memory_usage_summary} gives the values used
for the discussion below.

\begin{table}[htbp]
  \centering
  \caption{Average and peak memory usage by solver}
  \label{tab:memory_usage_summary}
  \begin{tabular}{lrr}
    \toprule
    Solver & Mean memory (KB) & Peak memory (KB) \\
    \midrule
    BruteForce & 6.43 & 11.07 \\
    A* h4: AC-3 pruned domain sum & 76.39 & 302.70 \\
    A* h1: empty cells & 92.03 & 485.05 \\
    A* h2: domain sum & 97.79 & 533.65 \\
    Forward chaining: data-driven fix-point & 285.23 & 603.54 \\
    Backtracking + forward chaining & 977.37 & 4734.27 \\
    Forward chaining \(\rightarrow\) backward chaining & 1200.26 & 6422.41 \\
    A* h3: min conflicts & 1537.23 & 20917.83 \\
    Backward chaining: generate-and-test SLD & 2106.31 & 7329.70 \\
    \bottomrule
  \end{tabular}
\end{table}

Memory mostly grows with board size, as expected. A larger Futoshiki board has
more cells, wider domains, and more inequality constraints to keep track of. The
growth is not perfectly smooth, though. Puzzle difficulty changes the amount of
work, and each solver stores a different kind of state while it runs. The hard
\(8 \times 8\) case makes this clear: A* h3 rises to about \(20\) MB, far above
the rest of the A* results.

Among the A* solvers, h1, h2, and h4 are the more memory-friendly choices. H1
and h2 are close together, averaging about \(92\) KB and \(98\) KB. H4 is even
steadier in this test: it averages about \(76\) KB and stays below about
\(303\) KB. That is a useful trade-off for AC-3 pruning. H4 keeps extra domain
information, but the pruning appears to prevent a much larger search frontier
later.

H3 is the exception. On smaller boards it looks harmless, sometimes close to the
other A* lines. The problem appears on harder larger puzzles. When the
min-conflicts heuristic does not keep the search focused, the open frontier
grows quickly, and memory follows. That is why h3 averages about \(1.5\) MB
even though many of its smaller cases are modest. Its peak, about \(20\) MB, is
the largest value in the whole memory plot.

The chaining solvers spend memory on logical state rather than search nodes.
Forward chaining keeps a fact base and updates it until inference has nothing
new to add. That makes it heavier than the light A* variants, but it stays
controlled: about \(285\) KB on average and about \(604\) KB at the peak. The
data-driven fix-point is the reason. The solver keeps inferred facts, applies
Modus Ponens while new facts are appearing, and stops when the fact set stops
changing. For a Futoshiki board, the useful value facts are tied to cells and
candidate values. Once those possibilities are exhausted, there is no growing
frontier left to store.

Backward chaining is heavier because it pays for grounding before it proves
anything. Requirement 2.3, Ground the KB, means the solver does not work
directly with the compact FOL schemas. For a \(9 \times 9\) board, those schemas
are expanded into a propositional knowledge base first. Some axioms create
\(O(N^3)\) ground literals or clauses, and the heavier pairwise rules can reach
\(O(N^4)\). Backward chaining then proves goals over that expanded KB, which
explains why its average is about \(2.1\) MB and why larger boards create memory
spikes. The forward-to-backward chaining solver follows the same pattern, but
with lower memory: about \(1.2\) MB on average and about \(6.3\) MB at the peak.

Backtracking with forward chaining lands between the moderate and heavy groups.
It averages about \(977\) KB and reaches about \(4.6\) MB on the hardest large
case. That fits the hybrid design: it keeps the forward-chaining propagation
state, then adds branching state when propagation alone cannot solve the puzzle.

Brute force is the smallest line, averaging only about \(6\) KB. That number
needs some care. Brute force is memory-light because it keeps almost no solver
state, but that does not make it practical. It pays for that simplicity with
time, so the low memory result should not be mistaken for good scalability.

Overall, the plot separates the solvers into a few clear groups. Brute force is
tiny but too slow to scale. A* h1, h2, and h4 stay memory-efficient, with h4
giving the most stable A* profile in this benchmark. Forward chaining costs more
memory because it maintains logical state, but it remains controlled. Backward
chaining, forward-to-backward chaining, and backtracking with forward chaining
are the heavier approaches. The main warning from the plot is h3: it can look
normal on easy puzzles, then become the most memory-hungry solver once the
search frontier expands.
